\documentclass[10pt,a4paper]{amsart}
\usepackage[margin=19mm]{geometry}
\usepackage{amsmath,amssymb,mathtools}
\usepackage[scr=rsfs,cal=euler]{mathalfa}
\usepackage{xfrac}
\usepackage[unicode,hidelinks,bookmarks=false]{hyperref}
\newcommand{\ie}{i.e.\ }
\newcommand{\cf}{cf.\ }
\newcommand{\resp}{respectively\ }
\newcommand{\Subsection}{\subsection}
\providecommand{\nobreakdash}{\nobreak\hbox{-}}
\setlength{\emergencystretch}{2em}
\multlinegap0pt
\usepackage{amssymb}
\usepackage{enumerate}
\newtheorem{theorem}{Theorem}
\title{On transposed Poisson conformal algebras}
\author{Lamei Yuan, Hao Fang}
\date{Source: arXiv:2603.14735v1 (CC BY 4.0)}
\begin{document}
\maketitle

\noindent\emph{Source and licence.} On transposed Poisson conformal algebras. Version arXiv:2603.14735v1 (CC BY 4.0), distributed by arXiv under the Creative Commons Attribution 4.0 International licence, http://creativecommons.org/licenses/by/4.0/, as recorded in the arXiv licence metadata for this exact version and shown on the abstract page https://arxiv.org/abs/2603.14735v1. Full licence notice: \texttt{licenses/arxiv-2603.14735-CC-BY-4.0.txt}.\par\smallskip

\noindent\textit{Editorial scope.} Counted unit: the article's theorem eqh (source0.tex lines 335--348). The article's complete original proof of it is delivered with it (source0.tex lines 349--519, one \texttt{proof} environment of 171 lines). No mathematics is written, repaired or re-derived: the statement and the proof are the article's own lines, in the article's own notation, and no retained line is substituted or re-flowed.  Retained uncounted as the source-local context the proof needs: lines 177--181; lines 273--277; lines 286--288.  Nothing was omitted from any retained span, so the package carries no omission record.  No retained line contains a citation, so the wrapper carries no bibliography and no bibliography entry.  No retained line contains a figure environment, an image-inclusion command or drawing code, so no image is delivered and the package declares no asset dependency.  Editorial formatting only, in addition: an AMS \texttt{amsart} preamble that restates the article's own declarations and macros used by the retained lines, namely the environments \texttt{theorem} on separate counters, together with the standard packages those lines need.  The retained environments print in this document as Theorem 1 in order of appearance, whereas the article prints the same items with the numbering of its own full document.  The complete native source of the article as distributed in the arXiv e-print for this exact version is bundled unchanged as \texttt{sources/196428/source0.tex}.
   \begin{align}
        &[(\partial a)_{\lambda}b]=-\lambda[a_{\lambda}b],~[a_{\lambda}(\partial b)]=(\partial+\lambda)[a_{\lambda}b], &&\text{(Conformal sesquilinearity)}\\
        &[a_{\lambda}b]=-[b_{-\lambda-\partial}a], &&\text{(Skew-symmetry)}\\
        &[a_{\lambda}[b_{\mu}c]]=[[a_{\lambda}b]_{\lambda+\mu}c]+[b_{\mu}[a_{\lambda}c]], \quad \forall ~a,b,c\in {A}.&&\text{(Jacobi identity)} \label{Jacobi}
   \end{align}

   \begin{align}\label{rule3}
     %a\circ_{\lambda}[b_{\mu}c]=p\left([(a\circ_{\lambda}b)_{\lambda+\mu}c]+[b_{\mu}(a\circ_{\lambda}c)]\right),	
     %\frac{1}{p}\cdot (a\circ_{\lambda}[b_{\mu}c])=[(a\circ_{\lambda}b)_{\lambda+\mu}c]+[b_{\mu}(a\circ_{\lambda}c)],	
      2 (a\circ_{\lambda}[b_{\mu}c])=[(a\circ_{\lambda}b)_{\lambda+\mu}c]+[b_{\mu}(a\circ_{\lambda}c)],\quad\forall~ a,b,c\in \mathcal{L}.	
   \end{align}%\emph{$p$-\textsc{tpca}}, $\mathfrak{tpca}(p)$

	\begin{align}\label{rule5}
		2([a_{\lambda}b]\circ_{\mu}c)=[a_{{\lambda}}(b\circ_{\mu-\lambda}c)]-[b_{\mu-\lambda}(a\circ_{\lambda}c)],\quad\forall~ a,b,c\in \mathcal{L}.
	\end{align}

\begin{theorem}
   Let $(\mathcal{L} ,\circ_\lambda,[{\cdot_ \lambda \cdot}])$ be a transposed Poisson conformal algebra. Then the following identities hold:
   \begin{align}
      &x\circ_{\lambda}[y_{\mu}z]+y\circ_{\mu}[z_{-\partial-\lambda}x]+z\circ_{-\partial-\lambda-\mu}[x_{\lambda}y]=0, \label{eqh}\\
      &[x_{\lambda}y]\circ_{\mu}z+[y_{\mu-\lambda}z]\circ_{-\partial-\lambda}x+[z_{-\partial-\lambda}x]\circ_{-\partial-\mu+\lambda}y=0, \label{eqw}\\
      &[(h\circ_{\lambda}[x_{\gamma}y])_{\lambda+\mu}z]+[(h\circ_{\lambda}[y_{\mu-\gamma}z])_{-\partial-\gamma}x]+[(h\circ_{\lambda}[z_{-\partial-\gamma}x])_{-\partial-\mu+\gamma}y]=0, \label{eqs}\\
      &[[x_{\gamma}y]_{\mu}(h\circ_{\lambda}z)]+[[y_{\mu-\gamma}z]_{-\partial-\lambda-\gamma}(h\circ_{\lambda}x)]+[[z_{-\partial-\gamma}x]_{-\partial-\lambda-\mu+\gamma}(h\circ_{\lambda}y)]=0, \label{equ}\\
      &[h_{\lambda}x]\circ_{\lambda+\gamma}[y_{\mu-\gamma}z]+[h_{\lambda}y]\circ_{\lambda+\mu-\gamma}[z_{-\partial-\gamma}x]+[h_{\lambda}z]\circ_{-\partial-\mu}[x_{\gamma}y]=0,  \label{eql}\\
      &[(u\circ_{\lambda}x)_{\mu}(v\circ_{\gamma}y)]+[(v\circ_{\gamma}x)_{\gamma+\mu-\lambda}(u\circ_{\lambda}y)]=2\left(u\circ_{\lambda}v\right)\circ_{\lambda+\gamma}[x_{\mu-\lambda}y], \label{ss1}\\
      &x\circ_{\mu-\lambda}[u_{\lambda}(v\circ_{\gamma}y)]+[(v\circ_{\gamma}x)_{\gamma+\mu-\lambda}u]\circ_{\gamma+\mu}y
      =\left(u\circ_{\lambda}v\right)\circ_{\lambda+\gamma}[x_{\mu-\lambda}y],\label{ss2}%~~p\neq1,
   \end{align}
   for all $x,y,z,h,u,v \in \mathcal{L}$.
\end{theorem}

\begin{proof}
Let $x,y,z,h,u,v \in \mathcal{L}$. In fact, Eq.~\eqref{eqw} is equivalent to Eq.~\eqref{eqh}. By Eqs.~\eqref{rule3} and \eqref{rule5}, we have
   \begin{align*}
      &{2} (x\circ_{\lambda}[y_{\mu}z])=[(x\circ_{\lambda}y)_{\lambda+\mu}z]+[y_{\mu}(x\circ_{\lambda}z)],	\\
      &{2} (y\circ_{\mu}[z_{-\partial-\lambda}x])=-{2} (y\circ_{\mu}[x_{\lambda}z])=-[(y\circ_{\mu}x)_{\lambda+\mu}z]-[x_{\lambda}(y\circ_{\mu}z)],\\
      &{2} (z\circ_{-\partial-\lambda-\mu}[x_{\lambda}y])={2} ([x_{\lambda}y]\circ_{\lambda+\mu}z)=[x_{{\lambda}}(y\circ_{\mu}z)]-[y_{\mu}(x\circ_{\lambda}z)].
   \end{align*}
Note that $[(x\circ_{\lambda}y)_{\lambda+\mu}z]=[(y\circ_{-\partial-\lambda}x)_{\lambda+\mu}z]=[(y\circ_{\mu}x)_{\lambda+\mu}z]$. Taking the sum of the three identities yields Eq.~\eqref{eqh}.
% Moreover, Eq.~\eqref{eqw} is equivalent to Eq.~\eqref{eqh}.

The proof of Eq.~\eqref{eqs} is more involved. First by Eq.~\eqref{rule3}, we have
   \begin{align*}
      2(h\circ_{\lambda}[[x_{\gamma}y]_{\mu}z])&=[(h\circ_{\lambda}[x_{\gamma}y])_{\lambda+\mu}z]+[[x_{\gamma}y]_{\mu}(h\circ_{\lambda}z)],\\
      -2(h\circ_{\lambda}[x_{\gamma}[y_{\mu-\gamma}z]])&=[(h\circ_{\lambda}[y_{\mu-\gamma}z])_{-\partial-\gamma}x]+[[y_{\mu-\gamma}z]_{-\partial-\lambda-\gamma}(h\circ_{\lambda}x)],\\
      2(h\circ_{\lambda}[y_{\mu-\gamma}[x_{\gamma}z]])&=[(h\circ_{\lambda}[z_{-\partial-\gamma}x])_{-\partial-\mu+\gamma}y]+[[z_{-\partial-\gamma}x]_{-\partial-\lambda-\mu+\gamma}(h\circ_{\lambda}y)].
   \end{align*}
Upon summing the three identities above and applying the Jacobi identity~\eqref{Jacobi}, we obtain
   \begin{eqnarray}\label{Id1}
      \begin{aligned}
      &[(h\circ_{\lambda}[x_{\gamma}y])_{\lambda+\mu}z]+[(h\circ_{\lambda}[y_{\mu-\gamma}z])_{-\partial-\gamma}x]+[(h\circ_{\lambda}[z_{-\partial-\gamma}x])_{-\partial-\mu+\gamma}y]+ \\
      &[[x_{\gamma}y]_{\mu}(h\circ_{\lambda}z)]+[[y_{\mu-\gamma}z]_{-\partial-\lambda-\gamma}(h\circ_{\lambda}x)]+[[z_{-\partial-\gamma}x]_{-\partial-\lambda-\mu+\gamma}(h\circ_{\lambda}y)]=0.
      \end{aligned}
   \end{eqnarray}
Then, applying the Jacobi identity to $[[x_{\gamma}y]_{\mu}(h\circ_{\lambda}z)]$ yields
%    \begin{align*}
%       [[x_{\gamma}y]_{\mu}(h\circ_{\lambda}z)]+[[y_{\mu-\gamma}(h\circ_{\lambda}z)]_{-\partial-\gamma}x]+[[(h\circ_{\lambda}z)_{-\partial-\gamma}x]_{-\partial-\mu+\gamma}y]=0
%    \end{align*}
   \begin{align*}
      [[x_{\gamma}y]_{\mu}(h\circ_{\lambda}z)]+[y_{\mu-\gamma}[x_{\gamma}(h\circ_{\lambda}z)]]=[x_{\gamma}[y_{\mu-\gamma}(h\circ_{\lambda}z)]]
   \end{align*}
and applying Eq.~\eqref{rule3} to $[y_{\mu-\gamma}(h\circ_{\lambda}z)]$ yields
   \begin{align*}
      [y_{\mu-\gamma}(h\circ_{\lambda}z)]=2\left(h\circ_{\lambda}[y_{\mu-\gamma}z]\right)-[(h\circ_{\lambda}y)_{\lambda+\mu-\gamma}z].
   \end{align*}
Hence, we arrive at
   \begin{align*}
      [[x_{\gamma}y]_{\mu}(h\circ_{\lambda}z)]+2[(h\circ_{\lambda}[y_{\mu-\gamma}z])_{-\partial-\gamma}x]=-[x_{\gamma}[(h\circ_{\lambda}y)_{\lambda+\mu-\gamma}z]]-[y_{\mu-\gamma}[x_{\gamma}(h\circ_{\lambda}z)]].
    %   [[x_{\gamma}y]_{\mu}(h\circ_{\lambda}z)]+p[(h\circ_{\lambda}&[y_{\mu-\gamma}z])_{-\partial-\gamma}x]\nonumber\\
    %   &=[[(h\circ_{\lambda}y)_{\lambda+\mu-\gamma}z]_{-\partial-\gamma}x]-[[(h\circ_{\lambda}z)_{-\partial-\gamma}x]_{-\partial-\mu+\gamma}y].
   \end{align*}
In a similar way, we have
%    \begin{align}
%       &[[y_{\mu-\gamma}z]_{-\partial-\lambda-\gamma}(h\circ_{\lambda}x)]+p[(h\circ_{\lambda}[z_{-\partial-\gamma}x])_{-\partial-\mu+\gamma}y]=[y_{\mu-\gamma}[x_{\gamma}(h\circ_{\lambda}z)]]-[[(h\circ_{\lambda}x)_{\lambda+\gamma}y]_{\lambda+\mu}z],\\
%       &[[z_{-\partial-\gamma}x]_{-\partial-\lambda-\mu+\gamma}(h\circ_{\lambda}y)]+p[(h\circ_{\lambda}[x_{-\partial-\mu+\gamma}y])_{\lambda+\mu}z]=[x_{\gamma}[(h\circ_{\lambda}y)_{\lambda+\mu-\gamma}z]]-[[y_{\mu-\gamma}(h\circ_{\lambda}x)]_{\lambda+\mu}z].
%    \end{align}
   \begin{align*}
      [[y_{\mu-\gamma}z]_{-\partial-\lambda-\gamma}(h\circ_{\lambda}x)]+2[(h\circ_{\lambda}&[z_{-\partial-\gamma}x])_{-\partial-\mu+\gamma}y] \nonumber\\
      &=[y_{\mu-\gamma}[x_{\gamma}(h\circ_{\lambda}z)]]-[[(h\circ_{\lambda}x)_{\lambda+\gamma}y]_{\lambda+\mu}z],\\
%       [[z_{-\partial-\gamma}x]_{-\partial-\lambda-\mu+\gamma}(h\circ_{\lambda}y)]+p[(h\circ_{\lambda}&[x_{-\partial-\mu+\gamma}y])_{\lambda+\mu}z] \nonumber\\
%       &=[x_{\gamma}[(h\circ_{\lambda}y)_{\lambda+\mu-\gamma}z]]-[[y_{\mu-\gamma}(h\circ_{\lambda}x)]_{\lambda+\mu}z].\label{yqq1}
%    \end{align}
% Next we turn to prove that
%    \begin{align*}
%       &[(h\circ_{\lambda}[x_{\gamma}y])_{\lambda+\mu}z]=[(h\circ_{\lambda}[x_{-\partial-\mu+\gamma}y])_{\lambda+\mu}z],\\
%       &[[(h\circ_{\lambda}x)_{\lambda+\gamma}y]_{\lambda+\mu}z]=-[[y_{\mu-\gamma}(h\circ_{\lambda}x)]_{\lambda+\mu}z].%=[[(h\circ_{\lambda}x)_{-\partial-\mu+\gamma}y]_{\lambda+\mu}z].
%    \end{align*}
% By the definition of conformal algebras, we compute directly that
%    \begin{align*}
%       &[(h\circ_{\lambda}[x_{-\partial-\mu+\gamma}y])_{\lambda+\mu}z]=\Big[\Big(h\circ_{\lambda}\Big(\sum_{n\in \mathbb{Z}_{+}}(-\partial-\mu+\gamma)^{(n)}(x_{[n]}y)\Big)\Big)_{\lambda+\mu}z\Big] \\
%       =&\Big[\Big(\sum_{n\in \mathbb{Z}_{+}}(-\partial-\lambda-\mu+\gamma)^{(n)}\Big(h\circ_{\lambda}(x_{[n]}y)\Big)\Big)_{\lambda+\mu}z\Big]
%       =\sum_{n\in \mathbb{Z}_{+}}\gamma^{(n)}\Big[\Big(h\circ_{\lambda}(x_{[n]}y)\Big)_{\lambda+\mu}z\Big] \\
%       =&\Big[\Big(h\circ_{\lambda}\Big(\sum_{n\in \mathbb{Z}_{+}}\gamma^{(n)}(x_{[n]}y)\Big)\Big)_{\lambda+\mu}z\Big]=[(h\circ_{\lambda}[x_{\gamma}y])_{\lambda+\mu}z],\\
% %    \end{align*}
% %    \begin{align*}
%       &-[[y_{\mu-\gamma}(h\circ_{\lambda}x)]_{\lambda+\mu}z]=[[(h\circ_{\lambda}x)_{-\partial-\mu+\gamma}y]_{\lambda+\mu}z]\\
%       =&\Big[\Big(\sum_{n\in \mathbb{Z}_{+}}(-\partial-\mu+\gamma)^{(n)}\Big((h\circ_{\lambda}x)_{[n]}y\Big)\Big)_{\lambda+\mu}z\Big]
%       =\sum_{n\in \mathbb{Z}_{+}}(\lambda+\gamma)^{(n)}\Big[\Big((h\circ_{\lambda}x)_{[n]}y\Big)_{\lambda+\mu}z\Big] \\
%       =&\Big[\Big(\sum_{n\in \mathbb{Z}_{+}}(\lambda+\gamma)^{(n)}\Big((h\circ_{\lambda}x)_{[n]}y\Big)\Big)_{\lambda+\mu}z\Big]
%       =[[(h\circ_{\lambda}x)_{\lambda+\gamma}y]_{\lambda+\mu}z].
%    \end{align*}
%     %   &[[(h\circ_{\lambda}x)_{-\partial-\mu+\gamma}y]_{\lambda+\mu}z]=\Big[\Big(\sum_{n\in \mathbb{Z}_{+}}(-\partial-\mu+\gamma)^{(n)}\Big((h\circ_{\lambda}x)_{[n]}y\Big)\Big)_{\lambda+\mu}z\Big] \\
%     %   =&\sum_{n\in \mathbb{Z}_{+}}(\lambda+\gamma)^{(n)}\Big[\Big((h\circ_{\lambda}x)_{[n]}y\Big)_{\lambda+\mu}z\Big]
%     %   =\Big[\Big(\sum_{n\in \mathbb{Z}_{+}}(\lambda+\gamma)^{(n)}\Big((h\circ_{\lambda}x)_{[n]}y\Big)\Big)_{\lambda+\mu}z\Big]\\
%     %   =&[[(h\circ_{\lambda}x)_{\lambda+\gamma}y]_{\lambda+\mu}z].
% Thus, Eq.~\eqref{yqq1} can be rewritten as
%    \begin{align}
      [[z_{-\partial-\gamma}x]_{-\partial-\lambda-\mu+\gamma}(h\circ_{\lambda}y)]+2[(h\circ_{\lambda}&[x_{\gamma}y])_{\lambda+\mu}z] \nonumber\\
      &=[x_{\gamma}[(h\circ_{\lambda}y)_{\lambda+\mu-\gamma}z]]+[[(h\circ_{\lambda}x)_{\lambda+\gamma}y]_{\lambda+\mu}z].
   \end{align*}
Upon summing the three identities above, we obtain
   \begin{eqnarray}
      \begin{aligned}\label{Id2}
      &2\left([(h\circ_{\lambda}[x_{\gamma}y])_{\lambda+\mu}z]+[(h\circ_{\lambda}[y_{\mu-\gamma}z])_{-\partial-\gamma}x]+[(h\circ_{\lambda}[z_{-\partial-\gamma}x])_{-\partial-\mu+\gamma}y]\right)+ \\
      &[[x_{\gamma}y]_{\mu}(h\circ_{\lambda}z)]+[[y_{\mu-\gamma}z]_{-\partial-\lambda-\gamma}(h\circ_{\lambda}x)]+[[z_{-\partial-\gamma}x]_{-\partial-\lambda-\mu+\gamma}(h\circ_{\lambda}y)]=0.
      \end{aligned}
   \end{eqnarray}
% By taking the difference between Eq.\eqref{Id1} and Eq.\eqref{Id2}, providing that $p\neq 1$, we arrive at Eq.~\eqref{eqs}.
Taking the difference between Eq.~\eqref{Id1} and Eq.~\eqref{Id2} yields Eq.~\eqref{eqs}.
Moreover, Eq.~\eqref{equ} follows by substituting Eq.~\eqref{eqs} into Eq.~\eqref{Id1}.

By Eq.~\eqref{rule3}, we obtain
   \begin{align}\label{yqq2}
      2\left([x_{\gamma}y]\circ_{\mu}[h_{\lambda}z]\right)=[([x_{\gamma}y]\circ_{\mu}h)_{\lambda+\mu}z]+[h_{\lambda}([x_{\gamma}y]\circ_{\mu}z)].
   \end{align}
Observe that
%    \begin{align*}
%       &[(h\circ_{\lambda}[x_{\gamma}y])_{\lambda+\mu}z]=[([x_{\gamma}y]\circ_{-\partial-\lambda}h)_{\lambda+\mu}z]
%       =\Big[\Big(\sum_{s \in \mathbb{Z}_{+}} (-\partial-\lambda)^{(s)}\left([x_{\gamma}y]_{(s)}h\right)\Big)_{\lambda+\mu}z\Big] \\
%       =&\sum_{s \in \mathbb{Z}_{+}} \mu^{(s)}\Big[\left([x_{\gamma}y]_{(s)}h\right)_{\lambda+\mu}z\Big]
%       =\Big[\Big(\sum_{s \in \mathbb{Z}_{+}} \mu^{(s)}\left([x_{\gamma}y]_{(s)}h\right)\Big)_{\lambda+\mu}z\Big]
%       =[([x_{\gamma}y]\circ_{\mu}h)_{\lambda+\mu}z].
%    \end{align*}
   \begin{align*}
      [(h\circ_{\lambda}[x_{\gamma}y])_{\lambda+\mu}z]=[([x_{\gamma}y]\circ_{-\partial-\lambda}h)_{\lambda+\mu}z]
      =[([x_{\gamma}y]\circ_{\mu}h)_{\lambda+\mu}z].
   \end{align*}
Thus, Eq.~\eqref{yqq2} can be rewritten as
   \begin{align*}
      2\left([h_{\lambda}z]\circ_{-\partial-\mu}[x_{\gamma}y]\right)=[(h\circ_{\lambda}[x_{\gamma}y])_{\lambda+\mu}z]+[h_{\lambda}(z\circ_{-\partial-\mu}[x_{\gamma}y])].
   \end{align*}
Similarly, by Eq.~\eqref{rule5}, we have
   \begin{align*}
      2([h_{\lambda}x]\circ_{\lambda+\gamma}[y_{\mu-\gamma}z])&=[(h\circ_{\lambda}[y_{\mu-\gamma}z])_{-\partial-\gamma}x]+[h_{{\lambda}}(x\circ_{\gamma}[y_{\mu-\gamma}z])],\\
      2([h_{\lambda}y]\circ_{\lambda+\mu-\gamma}[z_{-\partial-\gamma}x])&=[(h\circ_{\lambda}[z_{-\partial-\gamma}x])_{-\partial-\mu+\gamma}y]+[h_{{\lambda}}(y\circ_{\mu-\gamma}[z_{-\partial-\gamma}x])].
   \end{align*}
Upon summing the three identities above and applying Eqs.~\eqref{eqh} and \eqref{eqs}, we get Eq.~\eqref{eql}.

Finally, we turn to the proof of Eqs.~\eqref{ss1} and \eqref{ss2}. Since $(\mathcal{L} ,\circ_\lambda)$ is a {commutative associative conformal algebra}, we have
   \begin{align*}
        \left(u\circ_{\lambda}v\right)\circ_{\lambda+\gamma}[x_{\mu-\lambda}y]=u\circ_{\lambda}\left(v\circ_{\gamma}[x_{\mu-\lambda}y]\right).
   \end{align*}
By Eq.~\eqref{rule3}, we obtain
   \begin{align}\label{qq1}
      {4}(u\circ_{\lambda}\left(v\circ_{\gamma}[x_{\mu-\lambda}y]\right))=[(u&\circ_{\lambda}x)_{\mu}(v\circ_{\gamma}y)]+[(v\circ_{\gamma}x)_{\gamma+\mu-\lambda}(u\circ_{\lambda}y)] \nonumber\\
      &+[(u\circ_{\lambda}(v\circ_{\gamma}x))_{\gamma+\mu}y]+[x_{\mu-\lambda}(u\circ_{\lambda}(v\circ_{\gamma}y))].
   \end{align}
Observe that
   \begin{align}\label{qq2}
      {2}\left(u\circ_{\lambda}v\right)\circ_{\lambda+\gamma}[x_{\mu-\lambda}y]&=[((u\circ_{\lambda}v)\circ_{\lambda+\gamma}x)_{\mu+\gamma}y]+[x_{\mu-\lambda}((u\circ_{\lambda}v)\circ_{\lambda+\gamma}y)] \nonumber\\
      &=[(u\circ_{\lambda}(v\circ_{\gamma}x))_{\gamma+\mu}y]+[x_{\mu-\lambda}(u\circ_{\lambda}(v\circ_{\gamma}y))].
   \end{align}
Then Eq.~\eqref{ss1} follows from taking the difference between Eq.~\eqref{qq1} and Eq.~\eqref{qq2}.

Note that the first term on the left-hand side of Eq.~\eqref{ss1} can be rewritten as
%    \begin{align*}
%       &[(u\circ_{\lambda}x)_{\mu}(v\circ_{\gamma}y)]=[(x\circ_{-\partial-\lambda}u)_{\mu}(v\circ_{\gamma}y)]
%       =\Big[\Big(\sum_{s \in \mathbb{Z}_{+}} (-\partial-\lambda)^{(s)}(x_{(s)}u)\Big)_{\mu}(v\circ_{\gamma}y)\Big] \\
%       =&\sum_{s \in \mathbb{Z}_{+}} (\mu-\lambda)^{(s)}\left[(x_{(s)}u)_{\mu}(v\circ_{\gamma}y)\right]
%       =\Big[\Big(\sum_{s \in \mathbb{Z}_{+}} (\mu-\lambda)^{(s)}(x_{(s)}u)\Big)_{\mu}(v\circ_{\gamma}y)\Big]
%       =[(x\circ_{\mu-\lambda}u)_{\mu}(v\circ_{\gamma}y)]
%    \end{align*}
   \begin{align*}
      [(u\circ_{\lambda}x)_{\mu}(v\circ_{\gamma}y)]=[(x\circ_{-\partial-\lambda}u)_{\mu}(v\circ_{\gamma}y)]
      =[(x\circ_{\mu-\lambda}u)_{\mu}(v\circ_{\gamma}y)]
   \end{align*}
and by Eq.~\eqref{rule3}, we obtain
   \begin{align}\label{f5}
      [(x\circ_{\mu-\lambda}u)_{\mu}(v\circ_{\gamma}y)]={2}(x\circ_{\mu-\lambda}[u_{\lambda}(v\circ_{\gamma}y)])-[u_{\lambda}(x\circ_{\mu-\lambda}(v\circ_{\gamma}y))].
   \end{align}
For the second term, by Eq.~\eqref{rule5}, we get
   \begin{align}\label{f6}
      [(v\circ_{\gamma}x)_{{\gamma+\mu-\lambda}}(u\circ_{\lambda}y)]={2}([(v\circ_{\gamma}x)_{\gamma+\mu-\lambda}u]\circ_{\gamma+\mu}y)+[u_{\lambda}((v\circ_{\gamma}x)\circ_{\gamma+\mu-\lambda}y)].
   \end{align}
Since
%    \begin{align*}
%       x\circ_{\mu-\lambda}(v\circ_{\gamma}y)&=(x\circ_{\mu-\lambda}v)\circ_{\gamma+\mu-\lambda}y=(v\circ_{-\partial-\mu+\lambda}x)\circ_{\gamma+\mu-\lambda}y \\
%       &=\Big(\sum_{s \in \mathbb{Z}_{+}} (-\partial-\mu+\lambda)^{(s)}(v_{(s)}x)\Big)\circ_{\gamma+\mu-\lambda}y
%       =\sum_{s \in \mathbb{Z}_{+}} \gamma^{(s)}\left((v_{(s)}x)\circ_{\gamma+\mu-\lambda}y\right) \\
%       &=\Big(\sum_{s \in \mathbb{Z}_{+}} \gamma^{(s)}(v_{(s)}x)\Big)\circ_{\gamma+\mu-\lambda}y=(v\circ_{\gamma}x)\circ_{\gamma+\mu-\lambda}y,
%    \end{align*}
   \begin{align*}
      x\circ_{\mu-\lambda}(v\circ_{\gamma}y)&=(x\circ_{\mu-\lambda}v)\circ_{\gamma+\mu-\lambda}y=(v\circ_{-\partial-\mu+\lambda}x)\circ_{\gamma+\mu-\lambda}y
      =(v\circ_{\gamma}x)\circ_{\gamma+\mu-\lambda}y,
   \end{align*}
we have %$[u_{\lambda}(x\circ_{\mu-\lambda}(v\circ_{\gamma}y))]=[u_{\lambda}((v\circ_{\gamma}x)\circ_{\gamma+\mu-\lambda}y)]$.
   \begin{align*}
        [u_{\lambda}(x\circ_{\mu-\lambda}(v\circ_{\gamma}y))]=[u_{\lambda}((v\circ_{\gamma}x)\circ_{\gamma+\mu-\lambda}y)].
   \end{align*}
Hence, by taking the sum of Eqs.~\eqref{f5} and \eqref{f6} and applying Eq.~\eqref{ss1}, we arrive at Eq.~\eqref{ss2}.
% Combining Eqs.~\eqref{f5} and~\eqref{f6} with Eq.~\eqref{ss1}, we deduce Eq.~\eqref{ss2}.
\end{proof}

\end{document}
